\scanpage{17}\mspage{- 2 -}
\begingroup
\normalmarginpar
de plus, on vérifie fibre par fibre que
\[
\mathrm{R}^{1}\,\mathrm{p}_{*}\,\handwritten{\mathcal{O}}(\mathrm{ne})=0\quad\text{pour}\quad\mathrm{n}>0\qquad.
\]

Il en résulte, par Riemann-Roch, que pour $\mathrm{n}>0$, le faisceau $\mathrm{p}_{*}\handwritten{\mathcal{O}}(\mathrm{ne})$ est looalement libre de rang $\mathrm{n}$, et que sa formation commute à tout changement de base. La suite exacte de cohomologie fournit de plus une suite exacte
\[
0\handwritten{\longrightarrow}\mathrm{p}_{*}\handwritten{\mathcal{O}}(\mathrm{ne})
\handwritten{\longrightarrow}\mathrm{p}_{*}\handwritten{\mathcal{O}}((\mathrm{n}+1)\mathrm{e})
\handwritten{\longrightarrow}\omega^{\handwritten{\otimes}-(\mathrm{n}+1)}
\handwritten{\longrightarrow}0\quad(\mathrm{n}>0)\,.
\]

Par ailleurs, $\handwritten{\mathcal{O}}(\mathrm{e})$ étant localement libre de rang $1$, on vérifie fibre par fibre, compte tenu de (SGA 1 X 1.2), que
\[
\handwritten{\mathcal{O}}_{\mathrm{S}}\handwritten{\xrightarrow{\sim}}
\mathrm{p}_{*}\handwritten{\mathcal{O}}_{\mathrm{E}}\handwritten{\xrightarrow{\sim}}
\mathrm{p}_{*}\handwritten{\mathcal{O}}(\mathrm{e})\qquad.
\]

En d'autres termes, les $\mathrm{p}_{*}\handwritten{\mathcal{O}}(\mathrm{me})$ pour $1\leq\mathrm{m}\leq\mathrm{n}$ définissent une filtration de $\mathrm{p}_{*}\handwritten{\mathcal{O}}(\mathrm{ne})$, de gradué associé
\makeatletter\tagsleft@true\makeatother
\[
\mathrm{Gr}\;\mathrm{p}_{*}\handwritten{\mathcal{O}}(\mathrm{ne})
=\sum_{\substack{\mathrm{i}=0\\\mathrm{i}\handwritten{\neq}1}}^{\mathrm{n}}\omega^{\handwritten{\otimes}-\mathrm{i}}\qquad.
\tag*{(1.1)}
\]

Le morphisme $\mathrm{p}$ est d'intersection complète relative ; il en résulte que $\mathrm{Rp}^{!}\handwritten{\mathcal{O}}_{\mathrm{S}}\in\mathrm{D}^{+}(\mathrm{E})$ a un seul groupe de cohomologie non\margin{\handwritten{\unsure{réf|vérif|vif}}} nulle, $\mathrm{R}^{1}\,\mathrm{p}^{!}\handwritten{\mathcal{O}}_{\mathrm{S}}$, et que ce dernier est un faisceau inversible sur $\mathrm{E}$ ; on l'appelle le faisceau des différentielles régulières sur $\mathrm{E}$ et on le désigne par $\Omega^{\mathrm{reg}}_{\mathrm{E}/\mathrm{S}}$\margin{\handwritten{$\omega_{\mathrm{E}/\mathrm{S}}$}}. Là où $\mathrm{p}$ est lisse, on a $\Omega^{\mathrm{reg}}_{\mathrm{E}/\mathrm{S}}=\Omega^{1}_{\mathrm{E}/\mathrm{S}}$. La dualité montre que $\mathrm{R}^{1}\mathrm{p}_{*}\Omega^{\mathrm{rcg}}_{\mathrm{E}/\mathrm{S}}$ cst isomorphe à $\handwritten{\mathcal{O}}$, $\mathrm{p}_{*}\Omega^{\mathrm{rcg}}_{\mathrm{E}/\mathrm{S}}$ est donc \handwritten{/} localement libre, de formation compatible à tout changement\margin{\handwritten{réf}} de base et, raisonnant fibre par fibre, on en déduit que
% archive sheet number: 16
\endgroup
\clearpage
